\documentclass[11pt,a4paper]{article}
\usepackage[margin=1.05in]{geometry}
\usepackage{amsmath,amssymb,amsthm,mathtools}
\usepackage{booktabs,tabularx,enumitem,xcolor}
\usepackage{fontspec}
\setmonofont{DejaVu Sans Mono}[Scale=0.82]
\usepackage{microtype}
\usepackage[colorlinks=true,linkcolor=blue!50!black,urlcolor=blue!50!black,citecolor=blue!50!black]{hyperref}

\theoremstyle{plain}
\newtheorem{theorem}{Theorem}[section]
\newtheorem{proposition}[theorem]{Proposition}
\newtheorem{lemma}[theorem]{Lemma}
\newtheorem{corollary}[theorem]{Corollary}
\newtheorem*{theoremA}{Theorem A}
\theoremstyle{definition}
\newtheorem*{definition}{Definition}
\theoremstyle{remark}
\newtheorem*{remark}{Remark}

\newcommand{\ip}[2]{\langle #1,#2\rangle}
\newcommand{\hstar}{\mathfrak h^{\ast}}
\newcommand{\Ad}{\operatorname{Ad}}
\newcommand{\G}{\mathcal G}
\newcommand{\Sec}{\mathcal S}
\newcommand{\new}{_{\mathrm{new}}}
\newcommand{\out}{_{\mathrm{out}}}
\newcommand{\sw}{_{\mathrm{sw}}}
\newcommand{\lean}[1]{\texttt{#1}}

\title{The unstable curve at $q\out$ reaches the boundary of the star domain\\[4pt]
\large A proof without computation, for \S14.4 of \emph{Packings of Smoothed Polygons}}
\author{}
\date{November 2026}

\begin{document}
\maketitle

\noindent
This note is written for the authors of \emph{Packings of Smoothed Polygons} (T.~Hales and
K.~Vajjha). In \S14.4 the book computes the unstable curve $W^u(q\out)$ of the
Reinhardt-Poincar\'e map in Mathematica (Figure~14.2) and reads two facts off the computation:
along the curve the switching is cyclic, $3\to2\to1\to3\to\cdots$, and a trajectory starting on
the curve eventually reaches the boundary of the star domain before its next switch. The proof of
Theorem~14.4.1 rests on the second fact. The book remarks that the computation could be repeated
with interval arithmetic. This note proves both facts with no numerical computation. Every step is
checked in the Lean formalization of the book (\S9).

\section*{Summary: is the book's numerical argument correct?}

Yes. Both facts that the book reads off Figure~14.2 are true, and Theorem~A below proves them. An
independent computation in ball arithmetic, along 500 orbits of the unstable curve, also
reproduces the book's numbers. It places the threshold, beyond which a trajectory on the curve
exits before it switches, at $r^\ast\approx0.0649$ (the book says ``about $0.065$''), and the
curve of switching points ends at $(r,\tilde x_{21})\approx(0.2098,\,-3.0333)$ (the book gives
$(0.21,\,-3.03)$).

The argument's only weakness is that it was not certified. Figure~14.2 joins 434 computed points
by a piecewise linear curve, and, as the book notes, the computation was not repeated with
rigorous numerics. The surrounding proofs also have three gaps of exposition. The proof of Theorem~14.4.1 uses, without citing them,
the time reversal of the second part of Theorem~16.4.1 and Lemma~8.3.4. Theorem~14.3.3 states
without argument that near the fixed point the local stable manifold is the exceptional divisor.
Section~8 fills all three. None of this changes a conclusion of the book.

\section*{The idea of the proof}

We must show the following. Start at a point of the Poincar\'e section near $q\out$, off the
exceptional divisor ($r>0$), and iterate the Reinhardt-Poincar\'e map $F$: follow the control
$e_2$ until $e_1$ takes over, then rotate. The orbit keeps switching in the cyclic order, and
after finitely many steps an arc runs out of the star domain before it can switch.

A computation along the curve cannot settle this alone. Near the divisor the map stretches $r$ by
about $r_{\mathrm{scale}}\approx6.27$ per step, so an orbit that starts at radius $r$ needs about
$\log(1/r)/\log r_{\mathrm{scale}}$ steps to get away from the divisor. That number is unbounded,
so any computation must stop somewhere and hand over to a separate argument near the fixed point.
Each step also involves a switching time, the root of a transcendental equation.

The proof avoids both problems with four ingredients.

\begin{enumerate}[leftmargin=*,itemsep=3pt]
\item \emph{One arc is explicit} (\S2). With the control fixed at $e_2$, the state, the costate
and the switching functions have closed forms. The arc leaves the star domain at an explicit
time, and each switching function is a sum of four exponential terms.
\item \emph{A point of the section is a few numbers} (\S\S2--5). Six numbers, tied by two
relations, describe a point of the section, and each has a meaning: a distance to a side of the
star domain, the gap between two Hamiltonians, or the slope of a switching function. One step of
$F$ becomes an explicit rule for updating them. The rule involves one more number, $s$, which
records when the switch happens.
\item \emph{A budget} (\S5). The distance $\varepsilon=\rho_1$ drops by exactly $2r$ at each
step. If the radii of the switching points never decrease, an orbit starting at radius $r_0$ can
take at most $\varepsilon/(2r_0)$ steps.
\item \emph{A trapping region} (\S\S6--7). Five sign conditions on the six numbers define a
region $\G$ of the section. It contains every point near $q\out$ with $r>0$, and in it the radius
grows at each step. One step from $\G$ either stays in $\G$ or reaches an arc that leaves the
star domain. With the budget, every orbit from $\G$ ends with such an arc.
\end{enumerate}

The unstable curve is never located. The argument covers every point near $q\out$ with $r>0$,
on the curve or not.

Most of the work is to show that $\G$ is kept. Three of its five conditions carry over to the
next point by themselves; the other two need an argument, and it is short. The relation
$\det\Lambda_1=9/4$ can be rearranged so that $q>0$ is impossible while $p$ is small, and at the
next point $p$ is small because the arc crossed the switching wall before reaching the side of the
star domain. When the next point lies too far from the singular state for this, its arc runs into
a cusp of the star domain and leaves it.

\section{Setting and result}

\paragraph{Notation.} We write $\ip AB=\operatorname{tr}(AB)$, so that $\det A=-\tfrac12\ip AA$
for traceless $A$. For a traceless matrix $X$ with entries $x_{ij}$, the star functions of
Corollary~3.1.6 are
\[
\rho_0=\frac{x_{21}+\sqrt3\,x_{11}}{\sqrt3},\qquad
\rho_1=\frac{x_{21}-\sqrt3\,x_{11}}{\sqrt3},\qquad
\rho_2=-\frac{3x_{12}+x_{21}}{2\sqrt3}.
\]
The star domain $\hstar$ is the set where $\det X=1$ and all three are positive. Its three sides
are $\rho_i=0$, and on $\det X=1$ the star functions satisfy
$\rho_0\rho_1+\rho_1\rho_2+\rho_2\rho_0=1$. At the singular state $X=J$ all three equal
$1/\sqrt3$.

We write $Z_1,Z_2,Z_3$ for the control matrices $Z_{100},Z_{010},Z_{001}$ of the vertices
$e_1,e_2,e_3$. Each pairing with a control matrix is a star function:
\begin{equation}\tag{1.1}
\ip X{Z_1}=-\tfrac{2}{\sqrt3}\,\rho_2,\qquad
\ip X{Z_2}=-\tfrac{2}{\sqrt3}\,\rho_1,\qquad
\ip X{Z_3}=-\tfrac{2}{\sqrt3}\,\rho_0 .
\end{equation}
With $\lambda_{\mathrm{cost}}=-1$ and $\hat\Lambda=\Lambda_1+\tfrac32J$, the Hamiltonian of the
vertex $e_j$ is
\[
H_j=\ip{\hat\Lambda}X-\frac{\ip{\Lambda_R}{Z_j}}{\ip X{Z_j}},
\]
and $\chi_{ij}=H_i-H_j$. The switching function $\chi_{ij}$ has a pole on a side of $\hstar$, so we
clear the denominators:
\[
N_{ij}=\chi_{ij}\,\ip X{Z_i}\ip X{Z_j}
=\ip{\Lambda_R}{Z_j}\ip X{Z_i}-\ip{\Lambda_R}{Z_i}\ip X{Z_j}.
\]
Inside $\hstar$ the factor $\ip X{Z_i}\ip X{Z_j}$ is positive by (1.1), so $N_{ij}$ has the sign
of $\chi_{ij}$.

\paragraph{The section.} Let $\Sec$ be the set of triples $w=(X,\Lambda_1,\Lambda_R)$ with
$X\in\hstar$, $\Lambda_1$ and $\Lambda_R$ traceless, $\det\Lambda_1=9/4$, $\ip X{\Lambda_R}=0$ and
$H_2=H_3=0$. These are the states an extremal can have at the moment it switches from $e_3$ to
$e_2$: the two vertices tie, and the maximized Hamiltonian is $0$. The chart $(r,\tilde x)$ of
\S14.1 parametrizes $\Sec$, and the radius $r=x_{11}=(\rho_0-\rho_1)/2$ measures the distance
from the singular state.

\paragraph{The map.} From $w\in\Sec$ follow the control $e_2$. If $\chi_{21}$ vanishes for a first
time $t\sw>0$ while $X$ is still in $\hstar$, put $F(w)=\Ad_{R^{-1}}\,w(t\sw)$. At $t\sw$ the
vertex $e_1$ ties with $e_2$. The rotation relabels the vertices ($e_1\to e_2$, $e_2\to e_3$), so
$F(w)$ is again a point where an $e_3$-arc gives way to an $e_2$-arc. This is the
Reinhardt-Poincar\'e map of \S14.3; Lemma~14.3.1 extends it analytically to a neighbourhood of
the fixed point $(0,\tilde x\out)$.

\begin{theoremA}
There is a neighbourhood $U$ of $(0,\tilde x\out)$ with the following property. For every chart
point of $U$ with $r>0$, with section point $w_0$, there is $K\ge0$ such that the iterates
$w_j=F^j(w_0)$ satisfy:
\begin{enumerate}[label=\textup{(\arabic*)},itemsep=1pt]
\item for $j\le K$, the control after the switch at $w_j$ is $e_2$: $H_1<H_2$ at $w_j$, and
$H_3<H_2$ just after $w_j$ along the $e_3$-flow;
\item for $j<K$, along the $e_2$-arc from $w_j$ the vertex $e_2$ is the strict maximizer of
$H_1,H_2,H_3$ on $(0,t\sw)$, and $H_1>H_2$ at times arbitrarily close after $t\sw$;
\item the $e_2$-arc from $w_K$ reaches the boundary of $\hstar$ at some time $T$, with $e_2$ the
strict maximizer on $(0,T)$.
\end{enumerate}
\end{theoremA}

Items (1) and (2) say that the switching is cyclic. Item (3) says that the next step of $F$ is not
defined, because the trajectory leaves the star domain first. The proof of Theorem~14.4.1 needs
this for the points of the local unstable curve (\S8).

\section{One arc in closed form}

Everything below rests on the closed-form solutions of Chapter~7. Two features make the
$e_2$-arc simple: the pairing $\ip X{Z_2}$ stays constant along it, and the two matrix
exponentials in the solution grow at the same rate. A few numbers read off the starting point
then determine the whole arc.

\subsection*{The state}

For $X_0\in\hstar$ we put
\[
\varepsilon=\rho_1(X_0),\qquad Y=1+\frac{\rho_2(X_0)}{\varepsilon},\qquad
Y'=1+\frac{\rho_0(X_0)}{\varepsilon},\qquad a=Y-2,\qquad b=Y'-2 .
\]
Here $\varepsilon$ is the distance to the side $\rho_1=0$, which stays constant along the arc, and
$Y$ measures the distance $\rho_2$ to the side through which the arc will leave, relative to
$\varepsilon$. At the singular state $Y=Y'=2$, so $a$ and $b$ measure the displacement from it;
near the fixed point both are of order $r$. The relation among the star functions becomes
\begin{equation}\tag{2.1}
\varepsilon^2\,YY'=1+\varepsilon^2,
\end{equation}
so $(\varepsilon,Y)$ determine $Y'$, and $Y,Y'>1$. The radius is $r=\varepsilon b/2$.

\begin{proposition}[the state]
With the control $e_2$ the state is $X(t)=e^{tP}X_0e^{-tP}$ with $P=Z_2/\ip{Z_2}{X_0}$
(Chapter~7). With $\tau=t/\varepsilon$,
\[
\rho_1(t)=\varepsilon,\qquad \rho_2(t)=\varepsilon\,(Ye^{-\tau}-1),\qquad
\rho_0(t)=\varepsilon\,(Y'e^{\tau}-1).
\]
So the arc stays in $\hstar$ for $0\le t<T:=\varepsilon\log Y$, and it leaves through the side
$\rho_2=0$ at $T$, transversally, with $d\rho_2/dt=-1$ there.
\end{proposition}

In the upper half-plane the state slides along a Euclidean ray toward the cusp $1/\sqrt3$ of the
ideal triangle $\hstar$ and crosses the side $\rho_2=0$. The exit time has a closed form; on an
arc, the only time without one is the switching time.

\begin{proof}
A frame diagonalizes $P$. Let $S=\left(\begin{smallmatrix}1&1\\0&\sqrt3\end{smallmatrix}\right)$,
$D=\operatorname{diag}(1,-1)$, and let $E_+$, $E_-$ be the elementary matrices with a $1$ in
position $(1,2)$, respectively $(2,1)$. Write a traceless matrix as
$M=S(m_0D+m_+E_++m_-E_-)S^{-1}$, briefly $M=(m_0,m_+,m_-)$. Then
$\ip MN=2m_0n_0+m_+n_-+m_-n_+$, and direct multiplication gives
\begin{equation}\tag{2.2}
Z_1=\tfrac1{\sqrt3}(-1,0,1),\qquad Z_2=\tfrac1{\sqrt3}(1,0,0),\qquad
Z_3=\tfrac1{\sqrt3}(-1,-4,0),\qquad J=\tfrac1{\sqrt3}(-1,-4,1).
\end{equation}
By (1.1) and (2.2),
\[
\rho_1=-x_0,\qquad \rho_2=x_0-\tfrac12x_+,\qquad \rho_0=x_0+2x_- .
\]
So $X_0=(-\varepsilon,x_+,x_-)$ with $x_+=-2\varepsilon Y$ and $x_-=\varepsilon Y'/2$, and
$x_+x_-=-(1+\varepsilon^2)$ by (2.1). Also $P=(-1/(2\varepsilon),0,0)$, so $\Ad_{e^{tP}}$ fixes
$D$ and multiplies $E_\pm$ by $e^{\mp\tau}$.
\end{proof}

\subsection*{The switching functions}

The costate also evolves by matrix exponentials. By Chapter~7,
\[
\Lambda_1(t)=\Ad_{e^{tP}}\,\Ad_{e^{-t\Omega}}\,\Lambda_1(0),\qquad \Omega=X_0+P .
\]
A basis of eigenvectors of $\operatorname{ad}\Omega$ makes this explicit.

\begin{lemma}[an eigenbasis]
The nilpotent matrices
\[
V_+=\Bigl(x_+,\;-\frac{x_+^2}{\varepsilon},\;\varepsilon\Bigr),\qquad
V_-=\Bigl(x_-,\;\varepsilon,\;-\frac{x_-^2}{\varepsilon}\Bigr)
\]
satisfy $[\Omega,V_+]=-\varepsilon^{-1}V_+$ and $[\Omega,V_-]=\varepsilon^{-1}V_-$. In the basis
$\Omega,V_+,V_-$ the only nonzero pairings are
\[
\ip\Omega\Omega=\frac1{2\varepsilon^2},\qquad \ip{V_+}{V_-}=\frac1{\varepsilon^2}.
\]
\end{lemma}

\begin{proof}
A direct computation, using $x_+x_-=-(1+\varepsilon^2)$.
\end{proof}

The eigenvalues $\pm1/\varepsilon$ of $\operatorname{ad}\Omega$ are the rates of $\Ad_{e^{tP}}$
on $E_\pm$. This coincidence, which amounts to $\det\Omega=\det P$, produces a resonant term
$\tau e^{\mp\tau}$ below.

At a point $w\in\Sec$ we define three costate numbers,
\begin{equation}\tag{2.3}
f_0=\frac{\sqrt3}{8\varepsilon^3}\,N_{21}(w),\qquad
q=-\frac{\sqrt3}{6\varepsilon}\,\ip{\hat\Lambda}{V_+},\qquad
p=\frac{2}{\sqrt3\,\varepsilon}\,\ip{\hat\Lambda}{V_-},
\end{equation}
and the abbreviations
\[
B=Y-1+q,\qquad A=1+f_0-B,\qquad B'=Y'-1+p,\qquad A'=1-B' .
\]

\begin{theorem}[normal form]
Let $w\in\Sec$. Along the $e_2$-arc from $w$, with $\tau=t/\varepsilon$,
\begin{align*}
\frac{\sqrt3}{8\varepsilon^3}N_{21}(t)&=f(\tau):=-1+Ae^{-\tau}+Be^{-2\tau}+Y\tau e^{-\tau},\\
\frac{\sqrt3}{8\varepsilon^3}N_{23}(t)&=g(\tau):=-1+A'e^{\tau}+B'e^{2\tau}-Y'\tau e^{\tau}.
\end{align*}
Moreover
\begin{equation}\tag{2.4}
f_0=-\frac{\sqrt3}{3\varepsilon^3Y'}\,\ip{\hat\Lambda}{X_0}.
\end{equation}
\end{theorem}

So once the state is known, three numbers decide how $e_1$ and $e_3$ compete with $e_2$ along the
whole arc.
\begin{itemize}[itemsep=1pt]
\item $f_0=f(0)$ is the initial lead of $e_2$ over $e_1$. By (2.4) it is positive exactly when the
reduced Hamiltonian $\ip{\hat\Lambda}X$ is negative at $w$.
\item $p=g'(0)$ is the rate at which $e_2$ pulls ahead of $e_3$ right after $w$. (We have
$g(0)=0$, since $e_2$ and $e_3$ tie at $w$.)
\item $q$ shapes $f$ beyond its initial value. When $q\le0$, the lead of $e_2$ over $e_1$ first
grows and then shrinks, so it vanishes exactly once (Lemma~4.3).
\end{itemize}
Near the fixed point $f_0$ is of order $r^3$ and $q,p$ are of order $r^2$, the scalings of the
blow-up.

\begin{proof}[Proof sketch]
By Chapter~7, $\Lambda_R(t)$ is given by a quadrature. Put
$M=\Ad_{e^{-tP}}\hat\Lambda(t)$, $\tilde\Lambda=\Ad_{e^{-tP}}\Lambda_R(t)$ and
$c=\ip{\Lambda_R}P$, and let $\Theta_\pm=\tilde\Lambda_\pm-x_\pm c$ in frame coordinates. Since
$\ip{\Lambda_R}{Z_2}=c\,\ip X{Z_2}$, (2.2) gives
\[
N_{21}=-\tfrac{2\varepsilon}{3}\,e^{-\tau}\,\Theta_+,\qquad
N_{23}=\tfrac{8\varepsilon}{3}\,e^{\tau}\,\Theta_- .
\]
Using $x_+x_-=-(1+\varepsilon^2)$ and $H_2\equiv0$, the equations for $\Lambda_R$ and $c$ reduce
to
\[
\Theta_+'=-2\ip M{V_+},\qquad \Theta_-'=2\ip M{V_-}.
\]
By Lemma~2.2,
$\ip{M(t)}{V_+}=e^{-\tau}\ip{\Lambda_1(0)}{V_+}+\tfrac32\ip J{\Ad_{e^{tP}}V_+}$, which is explicit by
(2.2). Integrating gives the four terms. The constant term is $-1$ and the coefficient of
$\tau e^{-\tau}$ is $Y$. The other two coefficients are fixed by $f(0)=f_0$ and by the definition
of $q$, using $\ip J{V_+}=-\frac{4\varepsilon}{\sqrt3}(Y^2-Y+1)$. The computation for $N_{23}$ is
the same, with $\ip J{V_-}=\frac{\varepsilon}{\sqrt3}(Y'^2-Y'+1)$. On $\Sec$, $g(0)=0$ because
$\chi_{23}=0$.
\end{proof}

\section{The relation \texorpdfstring{$\det\Lambda_1=9/4$}{det Lambda1 = 9/4}}

The section is four-dimensional (the chart $(r,\tilde x)$), but $\varepsilon,Y,f_0,q,p$ are five
numbers, so one relation remains. It comes from $\det\Lambda_1=9/4$, and it is the main tool of
\S6.

\begin{proposition}
On $\Sec$,
\begin{equation}\tag{3.1}
\bigl[Y'(f_0-q)-Yp\bigr]^2-4pq+2(YY'-1)(bq+ap)+2Y'(3-ab)f_0=0 .
\end{equation}
\end{proposition}

Read (3.1) as a budget. Near the singular state $a$ and $b$ are small, so $3-ab>0$. When $q\le0$
and $f_0,p\ge0$, every term is then nonnegative except $2(YY'-1)bq$. A positive $q$ would have to
be balanced by the term $-4pq$, so it requires $p$ to be large compared with $b$. Lemma~6.1 makes
this precise.

\begin{proof}
The three numbers (2.3), (2.4) are pairings of $\hat\Lambda$ with $V_+,V_-,X_0$, so they determine
$\Lambda_1$. In the basis of Lemma~2.2,
\[
\Lambda_1=\alpha\,\Omega+\beta V_++\gamma V_-,
\]
and the pairings of that lemma give $\alpha=2\varepsilon^2\ip{\Lambda_1}\Omega$,
$\beta=\varepsilon^2\ip{\Lambda_1}{V_-}$ and $\gamma=\varepsilon^2\ip{\Lambda_1}{V_+}$. With
(2.3), (2.4) and the values of $\ip J{V_\pm}$,
\[
\alpha=\sqrt3\,\varepsilon\bigl(1+\varepsilon^2(Y+Y'A-YB')\bigr),\qquad
\beta=-\tfrac{\sqrt3}{2}\,\varepsilon^3(Y'^2-B'),\qquad
\gamma=2\sqrt3\,\varepsilon^3(Y^2-B).
\]
By Lemma~2.2 again, $\det\Lambda_1=-\tfrac12\ip{\Lambda_1}{\Lambda_1}=-(\alpha^2+4\beta\gamma)/(4\varepsilon^2)$.
Using $\varepsilon^2(YY'-1)=1$, the equation $\det\Lambda_1=9/4$ becomes
\begin{equation}\tag{3.2}
\bigl(3-ab+Y'(f_0-q)-Yp\bigr)^2+3(YY'-1)^2=4(Y^2-B)(Y'^2-B'),
\end{equation}
and expanding gives (3.1).
\end{proof}

\section{One arc: switch or exit}

Three things can happen along the $e_2$-arc from a point of the section. The vertex $e_1$ may
catch up with $e_2$ inside the star domain; this is a switch, and the orbit continues. The state
may reach the boundary of the star domain first; this is an exit. Or $e_3$ may come back above
$e_2$, which would break the cyclic order. This section shows that the signs of $f_0,q,p$ and of
one more number decide which of these happens. We start with a tool that is also used in \S5.

\begin{lemma}[Poisson brackets]
Write $L_uH_v$ for the derivative of $H_v$ along the $e_u$-flow, and $P_u=Z_u/\ip X{Z_u}$. Then
\[
L_uH_v=\ip{\hat\Lambda}{[P_u-P_v,X]}+\ip{\Lambda_R}{[P_u,P_v]}
-\bigl(\ip{\Lambda_R}{P_u}+\ip{\Lambda_R}{P_v}\bigr)\ip X{[P_u,P_v]} .
\]
\end{lemma}

The right side is antisymmetric in $u$ and $v$, so $L_uH_u=0$ and $L_uH_v=-L_vH_u$: this is the
antisymmetry of the Poisson bracket of the vertex Hamiltonians. In particular each $H_u$ is
constant along its own flow, and $H_2\equiv0$ along the $e_2$-arc from a point of $\Sec$.

The next lemma rules out the third outcome.

\begin{lemma}[$e_3$ never comes back]
If $p>0$ and $Y'\ge2$, then $g>0$ on $(0,\infty)$.
\end{lemma}

\begin{proof}
Put $\gamma(y)=g(\log y)/y=-1/y+A'+B'y-Y'\log y$. Then $\gamma(1)=0$ and
$y^2\gamma'(y)=B'y^2-Y'y+1$, which equals $p>0$ at $y=1$. For $y\ge1$ its derivative is
$2B'y-Y'\ge2B'-Y'=Y'-2+2p>0$. So $\gamma$ increases on $[1,\infty)$.
\end{proof}

To follow the competition with $e_1$ it helps to measure time by $y=e^\tau$, which runs from $1$
at the start of the arc to $Y$ at the exit, and to look at
\begin{equation}\tag{4.1}
\varphi(y):=y\,f(\log y)=f_0+\beta_Y(y)-q\,\frac{y-1}{y},\qquad
\beta_Y(y):=Y\log y-(y-1)-(Y-1)\,\frac{y-1}{y}.
\end{equation}
The function $\varphi$ has the sign of the lead of $e_2$ over $e_1$. Its part $\beta_Y$ does not
involve the costate: it is the lead that the geometry of the arc alone produces when $f_0=q=0$.

\begin{lemma}[$e_1$ catches up at most once]
Suppose $f_0>0$ and $q\le0$, and let $y_+$ be the larger root of $y^2-Yy+B$. Then $\varphi$
increases on $[1,y_+]$ and decreases on $[y_+,\infty)$ to $-\infty$. So $f$ has exactly one zero
$\tau\sw>0$, where it changes sign, and $s:=e^{\tau\sw}$ satisfies
\[
s>y_+\ge\max(1,\,Y-1),\qquad s^2-Ys+B>0 .
\]
The switch comes before the exit ($s<Y$) exactly when $\varphi(Y)<0$. If $\varphi(Y)\ge0$, then
$f>0$ on $(0,\log Y)$.
\end{lemma}

So the sign of $\varphi(Y)$ decides the race between the switch and the exit. When the switch
happens, it falls in the last stretch $Y-1<y<Y$ of the arc.

\begin{proof}
We have $\varphi'(y)=-(y^2-Yy+B)/y^2$. At $y=1$ the quadratic equals $q\le0$, so $1$ lies between
its roots, and $\varphi(1)=f_0>0$. Also $y_+=\tfrac12\bigl(Y+\sqrt{(Y-2)^2-4q}\bigr)\ge\tfrac12(Y+|Y-2|)$.
\end{proof}

\begin{proposition}[one arc]
Let $w\in\Sec$ with $r>0$, $f_0>0$ and $p>0$.
\begin{enumerate}[label=\textup{(\arabic*)},itemsep=1pt]
\item The control after $w$ is $e_2$ \textup{(}item \textup{(1)} of Theorem~A\textup{)}.
\item If $q\le0$ and $\varphi(Y)<0$, the arc from $w$ switches to $e_1$ at
$t\sw=\varepsilon\tau\sw<T$ \textup{(}item \textup{(2)}\textup{)}.
\item If $q\le0$ and $\varphi(Y)\ge0$, the arc from $w$ leaves $\hstar$ at $T=\varepsilon\log Y$
before any switch \textup{(}item \textup{(3)}\textup{)}.
\end{enumerate}
\end{proposition}

\begin{proof}
Inside $\hstar$, $\chi_{21}$ has the sign of $f$ and $\chi_{23}$ that of $g$. For (1): $H_1<H_2$ at
$w$ is $f_0>0$. By Lemma~4.1 the slope of $H_3-H_2$ along the $e_3$-flow at $w$ is
$L_3(H_3-H_2)=L_2H_3$, the slope of $-\chi_{23}$ along the $e_2$-flow. Since $\chi_{23}(w)=0$ and
$\frac{d}{dt}N_{23}(0)=\frac{8\varepsilon^2}{\sqrt3}\,p$, this slope is $-2\sqrt3\,p/(Y'-1)<0$.
Parts (2) and (3) follow from Lemmas~4.2 and~4.3 (note that $r>0$ means $Y'>2$) and
Proposition~2.1.
\end{proof}

\section{One step of the map}

We now follow a point through one switch and the rotation, and read off the numbers of the next
point. The update is explicit in terms of $s=e^{\tau\sw}$. Since $Y-1<s<Y$, the number $s$ records
where in the last stretch of the arc the switch happens.

\begin{theorem}[the map]
Let $w\in\Sec$ with $r>0$, $f_0>0$, $q\le0$, $p>0$ and $\varphi(Y)<0$. Then $F(w)\in\Sec$, and its
numbers are
\begin{equation}\tag{5.1}
\varepsilon\new=\frac{\varepsilon(Y-s)}{s},\qquad
Y'\new=\frac{Y}{Y-s},\qquad
Y\new=\frac{Y's^2-2s+Y}{Y-s},
\end{equation}
\begin{equation}\tag{5.2}
f_{0,\mathrm{new}}=\Bigl(\frac{s}{Y-s}\Bigr)^2g(\tau\sw),\qquad
p\new=-\Bigl(\frac{s}{Y-s}\Bigr)^2f'(\tau\sw)=\frac{s^2-Ys+B}{(Y-s)^2}.
\end{equation}
The radius of $F(w)$ is $r\new=(\varepsilon-\varepsilon\new)/2$.
\end{theorem}

The formulas have a direct meaning. The rotation turns the exit side $\rho_2=0$ into the side
$\rho_1=0$, so $\varepsilon\new$ is the distance to the exit side at the moment of the switch. A
switch close to the exit ($s$ near $Y$) gives a next point close to a side, with $\varepsilon\new$
small and $Y\new$, $Y'\new$ large. By (5.2), $f_{0,\mathrm{new}}$ is the lead of $e_2$ over $e_3$
at the switch and $p\new$ is the rate at which $e_1$ overtook $e_2$, both rescaled.

\begin{proof}
Since $\ip{\Ad_{R^{-1}}M}{Z_j}=\ip M{Z_{j-1}}$ (indices mod 3) and $\ip{\hat\Lambda}X$ is
$\Ad_R$-invariant, the rotation sends $H_j$ to $H_{j-1}$, $N_{ij}$ to $N_{i-1,j-1}$, and
$(\rho_0,\rho_1,\rho_2)$ to $(\rho_1,\rho_2,\rho_0)$. Along the arc $H_2\equiv0$, and $H_1=H_2$ at
$t\sw$, so the new $H_2$ and $H_3$ vanish. The quantities $\det\Lambda_1$ and $\ip X{\Lambda_R}$
are conserved and $\Ad$-invariant, and $X(t\sw)\in\hstar$. So $F(w)\in\Sec$. Proposition~2.1 at
$\tau\sw=\log s$ gives (5.1) and $r\new=(\rho_0-\rho_1)\new/2=(\varepsilon-\varepsilon\new)/2$.

For (5.2) we need $\Lambda_1(t\sw)$, since by (2.3) and (2.4) the numbers $f_{0,\mathrm{new}}$ and
$p\new$ are pairings of $\hat\Lambda$ with matrices built from the new state. The expansion
$\Lambda_1(0)=\alpha\Omega+\beta V_++\gamma V_-$ of \S3 can be carried along the arc. The map
$\Ad_{e^{-t\Omega}}$ multiplies $\Omega,V_+,V_-$ by $1,e^{\tau},e^{-\tau}$ (Lemma~2.2), and
$\Ad_{e^{tP}}$ sends $\Omega$ to $X(t)+P$ and $V_\pm$ to $e^{\pm\tau}V_\pm(X(t))$, where
$V_\pm(X)$ denotes the matrices of Lemma~2.2 built from $X$ in place of $X_0$. Hence
\begin{equation}\tag{5.4}
\Lambda_1(t)=\alpha\bigl(X(t)+P\bigr)+\beta e^{2\tau}\,V_+(X(t))+\gamma e^{-2\tau}\,V_-(X(t)) .
\end{equation}
Pairing (5.4) at $t=t\sw$, after the rotation, with the new state and with $V_-$ of the new state
gives $p\new=(s^2-Ys+B)/(Y-s)^2$ and an expression for $f_{0,\mathrm{new}}$ that is rational in
$s$. Eliminating $\tau\sw$ with $f(\tau\sw)=0$, which reads $Y\tau\sw=s-A-B/s$, turns these into
the two expressions in terms of $g(\tau\sw)$ and $f'(\tau\sw)$.
\end{proof}

Three of the conditions we need carry over to the next point by themselves, and the radius grows
by an explicit factor.

\begin{corollary}
Under the hypotheses of Theorem~5.1, $F(w)$ has $r>0$, $f_0>0$ and $p>0$, so the control after
$F(w)$ is $e_2$ \textup{(}Proposition~4.4\textup{)}. If $Y\ge2$, then
\begin{equation}\tag{5.3}
\frac{r(F(w))}{r(w)}>\frac{Y-2}{(Y-1)(Y'-2)} .
\end{equation}
\end{corollary}

\begin{proof}
$g(\tau\sw)>0$ by Lemma~4.2, and $s^2-Ys+B>0$ by Lemma~4.3. The radius $r\new$ is positive when
$\varepsilon\new<\varepsilon$, that is, when $s>Y/2$, which holds since $s>\max(1,Y-1)$. For
(5.3), write $r\new=\frac\varepsilon2\bigl(1-\frac{Y-s}{s}\bigr)$ and use $\frac{Y-s}{s}<\frac1{Y-1}$,
which holds because $s>Y-1$; also $r=\varepsilon(Y'-2)/2$.
\end{proof}

Near the fixed point the right side of (5.3) tends to about $2.57$, while the true stretching
factor is $r_{\mathrm{scale}}\approx6.27$. The bound is crude, but it is all we need. The last
ingredient of this section is the budget.

\begin{corollary}[the budget]
If $w_0,\ w_1=F(w_0),\ \dots,\ w_n$ are defined, then
\[
\sum_{j=1}^n r(w_j)=\tfrac12\bigl(\rho_1(w_0)-\rho_1(w_n)\bigr)<\tfrac12\,\rho_1(w_0).
\]
\end{corollary}

\begin{proof}
$\rho_1(w_{j+1})=\varepsilon\new=\rho_1(w_j)-2r(w_{j+1})$ by Theorem~5.1, and $\rho_1>0$.
\end{proof}

Each step pays for its radius out of $\varepsilon$. If the radii never decrease, the orbit stops
after at most $\rho_1(w_0)/(2r(w_0))$ steps. This replaces following the curve.

\section{The trapping region}

\begin{definition}
$\G$ is the set of $w\in\Sec$ that satisfy the five conditions
\[
r>0,\qquad f_0>0,\qquad p>0,\qquad q\le0,\qquad (Y-1)(Y'-2)\le Y-2 .
\]
\end{definition}

The first condition says that $w$ is off the exceptional divisor. The next two say that $e_1$
starts below $e_2$ and that $e_3$ falls behind $e_2$. The condition $q\le0$ makes $e_1$ catch up
at most once (Lemma~4.3). The last one, the \emph{expansion condition}, says that the radius
grows, by (5.3); with $r>0$ it forces $Y>2$ (so $a>0$) and $b<1$.

By Corollary~5.2, the first three conditions hold at the next point automatically. The whole
difficulty lies in the last two, and the plan is as follows. Lemma~6.1 turns the budget reading of
(3.1) into a sign test: when $p$ is small, $q$ cannot be positive. Lemma~6.2 shows that $p$ is
always small at the next point. Lemma~6.3 shows that the expansion condition carries over unless
the next point is far from the singular state ($Y\new\ge16/5$), and Lemma~6.4 shows that a point
that far out exits on its next arc.

\begin{lemma}[the relation as a sign test]
Let $w\in\Sec$ with $Y\ge2$, $Y'>2$, $f_0\ge0$, $p\ge0$ and
\begin{equation}\tag{6.1}
2p<(2Y'-1)(Y'-2).
\end{equation}
Then: \textup{(i)} if $(Y-2)(Y'-2)\le3$, then $q\le0$; \textup{(ii)} in any case
$f_0\ge q\,(1-1/Y)$.
\end{lemma}

Part (i) is the precise form of the budget: with $p$ small, the relation leaves no room for
$q>0$. Part (ii) also holds far from the singular state, and Lemma~6.4 uses it there.

\begin{proof}
Rearrange (3.1) as
\begin{equation}\tag{6.2}
2Y'(ab-3)\,f_0=2q\bigl[(YY'-1)b-2p\bigr]+\bigl[Y'(f_0-q)-Yp\bigr]^2+2(YY'-1)\,ap .
\end{equation}
Since $Y\ge2$, (6.1) gives $(YY'-1)b\ge(2Y'-1)b>2p$, so the bracket in the first term on the right
is positive. If $q\le0$, part (ii) is clear. Suppose $q>0$. Then the right side of (6.2) is
positive, so $ab>3$, which proves (i). Dropping the last two terms of (6.2) gives
$Y'(ab-3)f_0\ge q\bigl[(YY'-1)b-2p\bigr]$, and therefore
\[
YY'(ab-3)\Bigl[f_0-q\Bigl(1-\frac1Y\Bigr)\Bigr]\ \ge\ q\,D,\qquad
D=Y\bigl[(YY'-1)b-2p\bigr]-Y'(Y-1)(ab-3).
\]
With $a=Y-2$ one checks that
\[
D=Y\bigl[(2Y'-1)b-2p\bigr]+abY'+3Y'(Y-1),
\]
which is positive by (6.1). This proves (ii).
\end{proof}

Why is $p$ small at the next point? By (5.2), $p\new$ is the rescaled rate at which $e_1$
overtook $e_2$, and the overtaking happens in the last stretch of the arc, before the exit. There
the rate is limited by how far into that stretch the switch falls, the number $c=s-(Y-1)<1$. The
displacement $b\new$ of the next point has no such limit.

\begin{lemma}[$p$ is small at the next point]
Let $w\in\G$ with $\varphi(Y)<0$. Then $F(w)$ satisfies \textup{(6.1)}, and $Y\new>2$.
\end{lemma}

\begin{proof}
By (5.1),
\[
2Y'\new-1=\frac{Y+s}{Y-s},\qquad Y'\new-2=\frac{2s-Y}{Y-s}.
\]
Put $c=s-Y+1$; by Lemma~4.3, $0\le c<1$ and $s>1$. Since $q\le0$, the numerator of $p\new$
satisfies
\[
s^2-Ys+B=q+(s-1)c\ \le\ (s-1)c .
\]
Now $s-1\le2s-Y$ (their difference is $c$) and $2c<Y+s$, so
\[
2(s^2-Ys+B)\ \le\ 2(s-1)c\ \le\ 2c\,(2s-Y)\ <\ (2s-Y)(Y+s).
\]
Dividing by $(Y-s)^2$ gives (6.1) at $F(w)$. Finally $Y\new-2=(Y's^2-Y)/(Y-s)>0$, because
$Y's^2>2s^2\ge2(Y-1)^2\ge Y$ for $Y\ge2$.
\end{proof}

So Lemma~6.1 applies at the next point. Part (ii) always holds there, and part (i) gives
$q\new\le0$ provided the next point is not too far from the singular state. The expansion
condition behaves the same way.

\begin{lemma}[the expansion condition carries over unless $Y\new\ge16/5$]
Let $Y>2$, $Y'>2$ and $Y-1\le s<Y$, and define $Y\new$ and $Y'\new$ by \textup{(5.1)}. Then
\[
(Y-s)^2\bigl[(Y\new-2)-(Y\new-1)(Y'\new-2)\bigr]=\psi(s):=-3Y's^3+(2YY'+2)s^2-Y^2,
\]
and if $Y\new<16/5$, then $\psi(s)\ge0$: the expansion condition holds at the next point.
\end{lemma}

\begin{proof}
The identity is an expansion. Suppose $\psi(s)<0$ and $Y\new<16/5$; we derive a contradiction in
five steps.
\begin{enumerate}[itemsep=2pt]
\item $Y\new$ increases in $s$ on $[1,Y)$: its derivative is
$\bigl(Y's(2Y-s)-Y\bigr)/(Y-s)^2$, and $Y's(2Y-s)>YY's>Y$.
\item $\psi'(s)=s(4YY'+4-9Y's)$, so $\psi$ increases and then decreases on $(0,\infty)$; on any
interval its minimum is at an endpoint.
\item The function $\chi(s)=Y's^2+\tfrac65s-\tfrac{11}5Y=(Y-s)(Y\new-\tfrac{16}5)$ increases on
$(0,\infty)$, is negative at $0$ and positive at $Y$, so it has one root $s_\chi\in(0,Y)$. Since
$Y\new<16/5$, we have $Y-1\le s<s_\chi$.
\item At $s_\chi$ we have $Y\new=16/5$, and the expansion condition at the next point reads
$Y'\new\le28/11$, that is, $s_\chi\le\frac{17}{28}Y$, that is,
$\chi\bigl(\tfrac{17}{28}Y\bigr)=Y\bigl[\tfrac{289}{784}YY'-\tfrac{103}{70}\bigr]\ge0$, that is,
$YY'\ge\frac{5768}{1445}=3.9917\ldots$ This holds since $YY'>4$. So $\psi(s_\chi)\ge0$.
\item At $s=Y-1$, step~1 gives $Y\new(Y-1)=2+3a+2a^2+b(1+a)^2\le Y\new(s)<\frac{16}5$, which
forces $a<\frac13$. Then $\psi(Y-1)=a(2-a-2a^2)+b(1+a)^2(1-a)>0$.
\end{enumerate}
By step~2, $\psi\ge0$ on $[Y-1,s_\chi]$, which contains $s$. This is a contradiction.
\end{proof}

It remains to deal with next points far from the singular state. A large $Y\new$ means that the
switch came just before the exit, so the new arc starts close to a side and runs into a cusp of
$\hstar$. On such an arc the geometric part $\beta_Y$ of the switching function stays nonnegative
until the exit. A switch could then come only from the costate terms, and Lemma~6.1(ii) rules
that out.

\begin{lemma}[far points exit]
Let $w\in\Sec$ with $r>0$, $f_0>0$, $p>0$, $Y\ge16/5$ and $f_0\ge q(1-1/Y)$. Then the $e_2$-arc
from $w$ leaves $\hstar$ at $T=\varepsilon\log Y$, with $e_2$ the strict maximizer on $(0,T)$.
\end{lemma}

\begin{proof}
First, $G(Y):=\beta_Y(Y)=Y\log Y-(Y-1)-(Y-1)^2/Y$ is positive for $Y\ge16/5$. Its derivative,
$G'(Y)=\log Y-1+1/Y^2$, is positive there, since $\log Y\ge\log\frac{16}{5}>1$. At $Y=16/5$,
\[
G\bigl(\tfrac{16}5\bigr)=\tfrac{16}{5}\Bigl(\log\tfrac{16}{5}-\tfrac{297}{256}\Bigr),
\]
and
\[
\log\tfrac{16}{5}=2\operatorname{artanh}\tfrac{11}{21}
>2\Bigl(x+\frac{x^3}3+\frac{x^5}5+\frac{x^7}7\Bigr)\Big|_{x=11/21}>1.162>\tfrac{297}{256}
\]
since the series of $\operatorname{artanh}$ has positive terms. Next,
$\beta_Y'(y)=-(y-1)\bigl(y-(Y-1)\bigr)/y^2$, so $\beta_Y$ increases on $[1,Y-1]$ and decreases on
$[Y-1,Y]$. Since $\beta_Y(1)=0$ and $\beta_Y(Y)=G(Y)>0$, we get $\beta_Y\ge0$ on $[1,Y]$.

Now let $y\in(1,Y)$. If $q\le0$, then $\varphi(y)\ge f_0>0$ by (4.1). If $q>0$, then
\[
\varphi(y)\ \ge\ f_0-q\,\frac{y-1}{y}\ >\ f_0-q\Bigl(1-\frac1Y\Bigr)\ \ge\ 0 .
\]
So $f>0$ on $(0,\log Y)$. Lemma~4.2 gives $g>0$, and Proposition~2.1 gives the exit at $T$.
\end{proof}

The value $16/5$ matters only through $G(Y)>0$, which holds from $Y_\ast\approx3.163$ on; $16/5$
is a convenient rational number above $Y_\ast$.

\begin{theorem}[one step]
Let $w\in\G$ with $\varphi(Y)<0$. Then $F(w)$ is defined, the control after $F(w)$ is $e_2$,
$r(F(w))>r(w)$, and:
\begin{itemize}[itemsep=1pt]
\item if $Y\new\ge16/5$, the arc from $F(w)$ leaves $\hstar$ before switching;
\item if $Y\new<16/5$, then $F(w)\in\G$.
\end{itemize}
\end{theorem}

\begin{proof}
By Proposition~4.4, Theorem~5.1 and Corollary~5.2, $F(w)\in\Sec$ has $r>0$, $f_0>0$ and $p>0$,
and the control after it is $e_2$. By (5.3) and the expansion condition at $w$,
$r(F(w))>r(w)$. By Lemma~6.2, Lemma~6.1 applies at $F(w)$.

If $Y\new\ge16/5$: Lemma~6.1(ii) gives $f_{0,\mathrm{new}}\ge q\new(1-1/Y\new)$, and Lemma~6.4
gives the exit.

If $Y\new<16/5$: Lemma~6.3 gives the expansion condition at $F(w)$. That condition gives
$Y'\new-2<1$, and $Y\new-2<6/5$, so $(Y\new-2)(Y'\new-2)<3$. Lemma~6.1(i) gives $q\new\le0$,
and so $F(w)\in\G$.
\end{proof}

\section{Near the fixed point, and the proof of Theorem A}

The region $\G$ is useful only if the orbits we care about start in it. Near the divisor the five
conditions become, after rescaling by the natural powers of $r$, sign conditions on linear forms
in $\tilde x$. At the fixed point $\tilde x\out$ of the Fuller system these signs are right, with
room to spare.

\begin{lemma}
There is a neighbourhood $U$ of $(0,\tilde x\out)$ such that every chart point of $U$ with $r>0$
lies in $\G$.
\end{lemma}

\begin{proof}
Write $\tilde x=(\tilde x_{21},\tilde\ell_{11},\tilde\ell_{21})$. For small $r>0$ the chart point
lies in $\Sec$, since $X\to J$ and the star functions of $J$ equal $1/\sqrt3$. The quantities
\[
\frac{f_0}{r^3},\qquad \frac{q}{r^2},\qquad \frac{p}{r^2},\qquad \frac{Y-2}{r},\qquad
\frac{Y'-2}{r}
\]
are rational functions of $(r,\tilde x)$ whose denominators do not vanish at $r=0$. (The factor
$r^3$ in $f_0$ comes from $\det\Lambda_1=9/4$, which makes $\ip{\tilde\Lambda_1}J$ of order $r^2$.)
Their values at $r=0$ are
\[
-3(\tilde\ell_{11}+\tilde\ell_{21}\tilde x_{21}),\qquad
\sqrt3\,\tilde\ell_{11}+\tilde\ell_{21},\qquad
-2\tilde\ell_{21},\qquad
\sqrt3-3\tilde x_{21},\qquad 2\sqrt3,
\]
and the expansion condition, divided by $r$, tends to $-\sqrt3-3\tilde x_{21}$. Remark~14.2.1
locates the fixed point within $0.01$ of $\tilde x\out=(-2.39,\,-4.90,\,-1.12)$. There these
limits are about $6.63$, $-9.62$, $2.25$, $8.90$, $3.46$ and $5.44$, so the four linear forms have
the required signs, with margins of at least $1.1$. By continuity the five conditions hold near
$(0,\tilde x\out)$ for $r>0$.
\end{proof}

\begin{proof}[Proof of Theorem A]
Take $U$ from Lemma~7.1 and a chart point of $U$ with $r>0$, so that $w_0\in\G$. Item (1) holds
at $w_0$ by Proposition~4.4. Now suppose that $w_j\in\G$.
\begin{itemize}[itemsep=1pt]
\item If $\varphi(Y)\ge0$ at $w_j$, Proposition~4.4(3) gives item (3) with $K=j$.
\item If $\varphi(Y)<0$ at $w_j$, Proposition~4.4(2) gives item (2) at $w_j$. By Theorem~6.5, item
(1) holds at $w_{j+1}$, and either the arc from $w_{j+1}$ exits (item (3) with $K=j+1$) or
$w_{j+1}\in\G$ with $r(w_{j+1})>r(w_j)$.
\end{itemize}
While this continues the radii increase, so Corollary~5.3 gives $j\,r(w_0)<\rho_1(w_0)/2$. The
process therefore stops after fewer than $\rho_1(w_0)/(2r(w_0))$ steps, and it can stop only with
item (3).
\end{proof}

\begin{remark}
Theorem~A holds for every section point near the fixed point with $r>0$, so the location of the
unstable curve is never used.
\end{remark}

\begin{remark}
The two ways to exit match the book's description. If the switch would come after the exit, the
arc exits at once ($\varphi(Y)\ge0$); along the curve this is the regime beyond $r\approx0.065$.
Otherwise the switch may come just before the exit. The next point is then close to a side of
$\hstar$, with $Y\new$ large, and its arc runs into a cusp and exits (Lemma~6.4).
\end{remark}

\begin{remark}
A reference computation in ball arithmetic (python-flint/Arb), which the proof does not use,
follows 500 orbits of the curve. It finds the threshold $r^\ast\approx0.06495$ (the book says
about $0.065$; a comment in its \TeX\ source gives $0.0641$), and the curve of switching points
ends at $(r,\tilde x_{21})\approx(0.209847,\,-3.033319)$, against the book's $(0.21,\,-3.03)$.
From $r=10^{-6}$ on the curve the orbit has seven switching arcs; at the seventh $Y\new\approx38.2$,
and the next arc exits.
\end{remark}

\section{Theorem 14.4.1 and the local analysis of \S14.3}

\paragraph{Theorem 14.4.1.} The book's proof has four sentences. With Theorem~A they become a
complete proof.
\begin{enumerate}[itemsep=1pt]
\item The switching points of a trajectory that emanates from $q\out$ eventually lie on the local
unstable curve, hence in $U$ with $r>0$. This is the time reversal of the second part of
Theorem~16.4.1, which the book does not cite here.
\item A trajectory that returns to the singular locus must chatter into it. This is Lemma~8.3.4,
applied at the first return; the book does not cite it here either.
\item As long as the switching is cyclic, the switching points follow $F$. By Theorem~A the
trajectory reaches an arc that leaves the star domain, which contradicts the state constraint.
\item The second statement follows by time reversal.
\end{enumerate}

\paragraph{Theorem 14.3.3.} Its three claims are proved as follows.
\begin{itemize}[itemsep=1pt]
\item The fixed point is hyperbolic, by the three facts the book names: the Fuller-Poincar\'e map
contracts on the exceptional divisor, $r_{\mathrm{scale}}\approx6.27>1$, and the divisor is
invariant.
\item The local unstable curve is $C^\infty$ by Irwin's stable manifold theorem (Appendix~A.3),
which is proved in full in a separate Lean package.
\item The book states without argument that near the fixed point the local stable manifold is the
exceptional divisor. This is proved directly. Near the fixed point
$|F_1(r,\tilde x)|\ge|r|$, by the mean value theorem with $\partial_rF_1>1$ and $F_1(0,\cdot)=0$,
so no orbit with $r\ne0$ tends to the fixed point. And the divisor has a forward-invariant
neighbourhood of $\tilde x\out$ in which every orbit tends to it (linearized stability).
\end{itemize}
Remark~14.3.4 at $q_{\mathrm{in}}$ follows by time reversal, $F_{\mathrm{in}}=\tau F^{-1}\tau$.
The analytic extension of Lemma~14.3.1 is proved by the book's route, including the backward map
and analytic continuation.

\section{Where each step is checked}

The formalization\footnote{\url{https://github.com/vltanh/lean4-smoothed-polygon}} proves
Theorem~14.4.1 and Mahler's First conjecture using only the standard axioms of Lean:
\texttt{propext}, \texttt{Classical.choice} and \texttt{Quot.sound}. Its proof of the fact behind
Figure~14.2 is the proof of this note, step by step and with the same numbering. The table below
lists the main Lean declarations. They are in the folder \texttt{SmoothedPolygons/Ch14/} unless
stated otherwise, and the docstring of \texttt{ArcRegionInvariant.lean} lists every step.

\medskip
{\small
\noindent
\begin{tabularx}{\textwidth}{@{}l>{\raggedright\arraybackslash}X@{}}
\toprule
Step & Lean \\
\midrule
Theorem A (curve form) & \lean{UnstableCurveLeavesStarDomain}, proved as \lean{unstableCurveLeavesStarDomain} \\
Proposition 2.1 & \lean{ρ\_iso\_vfX\_one}, \lean{isStarState\_iso\_vfX\_one\_iff} (\lean{ArcFrame.lean}) \\
Lemma 2.2, Theorem 2.3 & \lean{commutator\_arcVe}, \lean{arc\_normalForm}, \lean{arcF0\_pos\_iff} \\
Proposition 3.1 & \lean{arcParams\_rel}, \lean{ArcParams.rel\_eq} \\
Lemma 4.1 & \lean{phaseHamSlope\_eq}, \lean{phaseHamSlope\_add\_swap} \\
Lemmas 4.2, 4.3 & \lean{ArcParams.g\_pos}, \lean{ArcParams.tauSw\_spec}, \lean{ArcParams.tauSw\_lt\_log\_iff} \\
Proposition 4.4 & \lean{switchEntersE2\_of\_arc}, \lean{arcSwitchesToE1\_of\_arc}, \lean{arcLeavesStarDomain\_of\_arc} \\
Theorem 5.1, (5.4) & \lean{arcParams\_reinPoincare}, \lean{ArcParams.step\_closed}, \lean{iso\_vfΛ₁\_one\_eq\_arcParams} \\
Corollaries 5.2, 5.3 & \lean{arc\_propagation}, \lean{sum\_radius\_reinPoincare\_iterate\_lt} \\
The region $\G$ & \lean{arcRegion} \\
Lemma 6.1 & \lean{ArcParams.q\_nonpos\_of\_relT}, \lean{ArcParams.f0\_ge\_of\_relT} \\
Lemma 6.2 & \lean{ArcParams.pNew\_lt}, \lean{ArcParams.two\_lt\_YNew} \\
Lemma 6.3 & \lean{ArcParams.E\_new\_iff}, \lean{ArcParams.YNew\_ge} \\
Lemma 6.4 & \lean{ArcParams.f\_pos\_of\_Y\_ge}, \lean{arcLeavesStarDomain\_of\_Y\_ge}; \lean{Real.log\_sixteen\_div\_five\_gt} \\
Theorem 6.5 & \lean{arcRegionInvariant\_sharp} \\
Lemma 7.1 & \lean{exists\_nhds\_lieChart\_mem\_arcRegion} \\
Proof of Theorem A & \lean{ArcRegionInvariant.reachesStarBoundary} \\
Theorem 14.4.1 & \lean{theorem14\_4\_1\_unconditional} (\lean{Ch16/MahlerFirstArcRegion.lean}), \lean{theorem14\_4\_1} \\
Theorem 14.3.3, A.3.1 & \lean{theorem14\_3\_3}, \lean{remark14\_3\_4}, \lean{theoremA\_3\_1} \\
\bottomrule
\end{tabularx}}

\medskip
The relation of Proposition~3.1 is checked in Lean by a direct computation in the frame
(\lean{arcParams\_rel\_identity}); the eigenbasis expansion in the proof above is the same
computation, organized to explain where the relation comes from.

\end{document}
